\documentclass[11pt,a4paper]{article}
\usepackage{ctex}
\usepackage{geometry}
\usepackage{amsmath,amssymb}
\usepackage{graphicx}
\usepackage{booktabs}
\usepackage{enumitem}
\usepackage{hyperref}
\geometry{left=2.5cm,right=2.5cm,top=2.5cm,bottom=2.5cm}
\hypersetup{colorlinks=true,linkcolor=blue}

\begin{document}
\title{仿生非轮式越障爬行机器人竞速赛实验方案}
\author{工程训练（制造工艺实习）小组}
\date{}
\maketitle

\section{实验项目概述}
本实验属于工程训练综合实训核心项目，小组自主完成一台仿生非轮式越障四足爬行机器人的完整开发，参与标准化赛道竞速竞赛。项目禁止轮式、履带行走机构；硬件限定使用ESP32-S3主控与8路MG90S舵机，不额外申请拓展传感器。可选用全自动自主运行模式或遥控模式参赛。

本小组选定\textbf{标准四足爬行机器人}作为唯一参赛方案。学生完整经历需求分析、方案选型、三维建模、数字化加工、机械装配、硬件接线、嵌入式步态编程、赛道调试、竞赛答辩全工程流程，锻炼机械设计、机电一体化集成、嵌入式软件开发与团队工程协作能力。

\section{实验目标}
\begin{enumerate}[itemsep=0.3em]
    \item 掌握仿生四足爬行机器人机械结构设计方法，完成非标零件自主建模与数字化加工；
    \item 掌握ESP32-S3嵌入式平台使用，实现8路舵机驱动与三角交替步态时序控制；
    \item 完成四足机器人整机机电联调，实现爬坡、隧洞通行、沟壑跨越等连续越障动作；
    \item 分析步态时序、机身结构参数对机器人行进速度、稳定性、越障能力的影响；
    \item 完成整机技术总结，梳理结构、电控、控制算法优化思路并开展答辩汇报。
\end{enumerate}

\section{竞赛赛道与技术约束}
\subsection{赛道参数}
赛道总长4m，一体式实训台，分为三段连续障碍：
\begin{enumerate}
    \item 组合式坡度越障区：40°上坡 + 45°下坡；全自动机器人必须翻越45°下坡段方认定有效全自动成绩；
    \item 狭窄矮洞隧洞区：通行区域尺寸长310mm，宽310mm，高度250mm；
    \item 弧形沟壑越障区：下凹曲面沟槽，检验自适应步态。
\end{enumerate}
机器人掉落、冲出赛道即刻终止本轮测试。

\subsection{机器人硬性约束}
\begin{enumerate}
    \item 整机外形上限：长$\le$300mm，宽$\le$300mm，高$\le$250mm；
    \item 禁止车轮、履带等轮式移动机构；
    \item 所有非标结构件自主设计加工，严禁成品机械套件；
    \item 关键受力构件不得使用泡沫、纸板等低强度材料；
    \item 硬件限制：仅使用课程标配ESP32-S3 + 8个MG90S舵机，不申请额外舵机与拓展传感器。
\end{enumerate}

\section{仿生四足爬行机器人方案说明}
\textbf{舵机分配：}四条机械腿，单腿2自由度（髋关节+膝关节），合计8台MG90S舵机。
\textbf{运动模式：}交替三角步态直行、原地转向、低速爬坡。
\begin{itemize}
    \item 优势：开源步态资料充足，零件数量少，装配简单，平地、坡道、沟壑场景均可适配；零基础短周期内可稳定完工，项目风险可控。
    \item 赛道适配：三角多点支撑保障坡道稳定；机身采用扁平设计满足矮洞限高；四肢交替支撑，平稳通过弧形沟壑。
    \item 优化措施：腿部加装防滑结构；优化重心高度，降低高速爬坡侧翻风险。
\end{itemize}

\section{实验硬件清单}
\begin{enumerate}
    \item 主控：ESP32-S3开发板；
    \item 执行机构：8台MG90S金属齿轮舵机；
    \item 辅助硬件：舵机独立电源模块、杜邦线、紧固件、焊接耗材；
    \item 遥控模式自备：遥控器、无线通信模块（课程不提供、不报销）。
\end{enumerate}

\section{实验实施步骤（7次实训任务）}
\subsection{任务1：项目导入、组队与方案确定}
学习竞赛完整规则，3～4人组建小组，确定分工；领取基础物料；确定采用标准四足爬行机器人方案，完成机身外形尺寸规划。

\subsection{任务2：整机方案设计与三维建模}
完成四足机器人整机三维建模，校核外形尺寸满足约束；规划舵机布局、供电方案；输出零件加工图纸、电气原理图初稿。

\subsection{任务3：结构零件数字化加工}
开展结构件3D打印、亚克力激光切割；零件去毛刺、尺寸校验；开展机械预装配，排查运动干涉。

\subsection{任务4：整机装配与基础步态调试}
完成机械装配与电控接线；统一完成舵机零点校准；编写三角交替步态程序，实现机器人直行、原地转向基础动作。

\subsection{任务5：越障步态专项调试}
分段测试坡道、隧洞、沟壑通行效果；迭代优化步态时序与机身重心控制；搭建全自动感知-决策-执行程序框架。

\subsection{任务6：完整赛道模拟联调}
多次全赛道试运行，记录故障卡点；迭代优化机械结构与控制代码；筹备答辩PPT，梳理整机结构、电控、步态优化内容。

\subsection{任务7：正式竞速竞赛 + 项目答辩}
参与两轮正式竞速测试；小组依次进行PPT答辩；提交全套工程图纸；清点、回收所有硬件物料。

\section{竞速评分计算方法}
总竞技得分 = 基础行进得分 + 竞速效率得分
\begin{enumerate}
    \item 基础行进得分（满分40）：
    $$S_{\text{基础}} = 40 \times \frac{L_{\text{实际}}}{L_{\text{总}}},\quad L_{\text{总}}=4\,\text{m}$$
    遥控模式原始基础得分$\times 0.7$折算。
    \item 竞速效率得分（满分60），按整体排名分四档：
    前25\%：60；25\%～50\%：57；50\%～75\%：48；后25\%：42。
\end{enumerate}

排名优先级：全自动模式 > 遥控模式；完赛队伍 > 未完赛队伍；同条件下耗时/行进距离排序。

\section{学习资源仓库}
电控教程、ESP32-S3示例代码、硬件原理图、接线图纸、3D参考模型存放于Gitee仓库：
\begin{center}
\url{shturl.cc/xHgiGnUs5MFMLUal4NnZi9cU9qTO0skfyIZhr}
\end{center}

\section{安全注意事项}
\begin{enumerate}
    \item 激光切割机、3D打印机、电烙铁必须遵守安全操作规程，防止烫伤、机械伤害；
    \item 舵机通电调试，禁止肢体靠近运动关节，避免夹伤；
    \item 锂电池规范充电，远离易燃物；
    \item 设备出现异响、过热、冒烟，立即断电并报告指导教师。
\end{enumerate}

\section{预期实验成果交付}
\begin{enumerate}
    \item 四足机器人零件二维加工图纸（DXF）；
    \item 整机三维模型源文件；
    \item 电气接线原理图；
    \item ESP32嵌入式四足三角步态完整源代码；
    \item 项目答辩PPT（包含四足机器人结构、电控、步态优化分析）；
    \item 小组实训技术总结报告。
\end{enumerate}

\end{document}